\documentclass[10pt,twocolumn,a4paper]{article}

\usepackage[utf8]{inputenc}
\usepackage[T1]{fontenc}
\usepackage{mathptmx}
\usepackage[scaled=0.92]{helvet}
\usepackage{courier}
\usepackage[a4paper,top=1.9cm,bottom=2.2cm,left=1.7cm,right=1.7cm,columnsep=0.65cm]{geometry}
\usepackage{titlesec}
\usepackage{booktabs}
\usepackage{caption}
\usepackage{enumitem}
\usepackage{xcolor}
\usepackage{url}
\usepackage[hidelinks]{hyperref}

\captionsetup[table]{font=small,labelfont=bf,skip=4pt}
\captionsetup[figure]{font=small,labelfont=bf,skip=4pt}

\titleformat{\section}{\normalfont\large\bfseries}{\thesection.}{0.6em}{}
\titleformat{\subsection}{\normalfont\normalsize\bfseries}{\thesubsection.}{0.6em}{}
\titlespacing*{\section}{0pt}{10pt}{4pt}
\titlespacing*{\subsection}{0pt}{8pt}{3pt}

\setlist[itemize]{leftmargin=1.1em,itemsep=1pt,topsep=3pt}
\setlist[enumerate]{leftmargin=1.4em,itemsep=1pt,topsep=3pt}

\renewcommand{\baselinestretch}{1.02}

\begin{document}

\twocolumn[
\begin{@twocolumnfalse}
\begin{center}
{\LARGE\bfseries BlueProof: Community-Verified Mangrove Restoration Monitoring\\[2pt]
with Vision-Language Models and Conditional Mobile Payments\par}
\vspace{10pt}
{\large Rooney Odhiambo\par}
\vspace{3pt}
{\normalsize GovTech Builders KE\\
Technical University of Mombasa, Kenya\\
\texttt{rooneyodhiambo2004@gmail.com}\par}
\vspace{10pt}
\end{center}

\begin{quote}
\small
\noindent\textbf{Abstract}---Kenya holds roughly 61{,}271 hectares of mangrove
forest, storing some of the densest carbon stocks of any tropical ecosystem, yet
the forest continues to be lost to cutting, urban encroachment and pest damage.
Community-led blue carbon projects demonstrably work, but they have not scaled:
the pioneering Kenyan project has issued only tens of thousands of credits in
fifteen years. We argue that the binding constraint is not a shortage of willing
communities but the cost and latency of monitoring, reporting and verification
(MRV). This paper presents BlueProof, a system design in which community
monitors submit geotagged plot photographs, a vision-language model (VLM)
performs bounded corroboration of those submissions, deterministic satellite
analysis provides canopy-scale context, and verified submissions trigger
immediate conditional mobile-money payment. We make the architectural role of
the VLM deliberately narrow. Recent evidence shows general-purpose VLMs degrade
sharply on field imagery relative to curated photographs and fabricate species
names at non-trivial rates, so BlueProof does not treat the model as an
authoritative taxonomic oracle. Instead the model corroborates a
monitor-declared species, assesses image legibility and plot condition, flags
cutting and pest damage, and abstains into human review when confidence is low.
We further argue for an incentive design that pays for verified \emph{reporting}
rather than for seedling survival, on the grounds that paying for outcomes
creates a direct incentive to conceal mortality. We situate the design against
the closest existing deployment, Ericsson's Connected Mangroves in Malaysia,
and argue that its premise transfers to Kenya while its apparatus does not:
fixed environmental sensors address growing conditions under corporate
sponsorship and volunteer labour, whereas the Kenyan constraint is attestation
cost in the absence of a sponsor, with mobile money affording a payment loop the
Malaysian model had no reason to build. We describe a reference
implementation and, because no field data has yet been collected, we set out a
pre-registered evaluation protocol with explicit metrics, baselines and
falsification criteria rather than reporting results. We state the ecological,
legal and equity limitations of the approach, including the risk that cheaper
verification expands credit supply without improving what a community monitor
earns.
\vspace{6pt}

\noindent\textbf{Keywords}---mangrove restoration, blue carbon, MRV,
vision-language models, participatory monitoring, mobile money, Kenya
\end{quote}
\vspace{12pt}
\end{@twocolumnfalse}
]

\section{Introduction}

Mangroves occupy a small fraction of the world's forested area but store carbon
at densities far above most terrestrial forests, buffer coastlines against
storm surge, and support inshore fisheries. In Kenya they cover approximately
61{,}271 hectares, of which about 35{,}678 hectares, some 62 percent, lie in
Lamu County~\cite{kairo2021}. Measured total ecosystem carbon stocks in Lamu
average $560.22 \pm 79.79$~Mg~C~ha$^{-1}$, with 70.3 percent of that carbon
held below ground in soil rather than in standing
biomass~\cite{kairo2021}. That distribution matters for monitoring: most of the
carbon at risk is invisible from above, and is released by disturbance of the
sediment rather than only by removal of the canopy.

The forest is under sustained pressure. Between 1990 and 2019 Lamu alone lost
1{,}739 hectares, a decline of roughly 60 hectares per year, with associated
emissions estimated at 30{,}840~Mg~CO$_2$e~yr$^{-1}$~\cite{kairo2021}. Drivers
differ by county. In Lamu and Kilifi the historical driver has been
overharvesting for construction poles and fuelwood. Around Mombasa, and
specifically in Tudor Creek and Port Reitz, urban encroachment and effluent are
dominant. Pest damage to \emph{Sonneratia alba}, a frontline species in the
seaward zone, has emerged as a compounding threat, while \emph{Xylocarpus
granatum} is selectively removed for carving and construction.

Community-led restoration and protection is not an untested idea in this
setting. Mikoko Pamoja at Gazi Bay, Kwale County, began in 2010 and is generally
described as the world's first community-led blue carbon project, involving
1{,}081 participating households~\cite{planvivo}. It works. What it has not done
is scale: the project has issued on the order of $2.2 \times 10^{4}$ certified
credits across roughly fifteen years~\cite{planvivo}. Against a national estate
of 61{,}271 hectares and a national loss rate measured in tens of hectares per
year, that is a rounding error.

The natural question is why a demonstrably functional model has not replicated.
Our position, which this paper formalises into a system design and an
evaluation protocol rather than asserting as established fact, is that the
binding constraint is the cost, latency and expertise-dependence of MRV. A
recent systematic treatment of MRV in carbon farming finds that monitoring
consumes the largest share of total MRV expenditure, that field sampling is the
most labour-intensive component, and that cost burdens fall disproportionately
on smaller operations to the point of threatening their economic
viability~\cite{vanderheyden2026}. The same work finds that digitalisation
measures such as automated data entry can halve data-entry time, and concludes
that policy support is likely required for MRV systems to be scalable for
smaller participants~\cite{vanderheyden2026}. A community forest association
managing tens of hectares of creek-side mangrove faces precisely this structure:
the verification machinery costs roughly what it costs regardless of project
size, so small projects are priced out of the market that would pay them.

\subsection{Contributions}

This paper makes four contributions.

\begin{enumerate}[label=(\arabic*)]
\item A system architecture for community mangrove MRV that separates
deterministic computation (vegetation indices, thresholds, payment rules,
ledger integrity) from probabilistic perception (image interpretation, natural
language), and makes that boundary explicit and auditable.
\item A deliberately \emph{bounded} role for a vision-language model, motivated
by published evidence on VLM failure modes in field imagery, in which the model
corroborates rather than determines species, and abstains into human review
under a calibrated confidence threshold.
\item An incentive design that pays community monitors for verified
\emph{reporting} rather than for seedling survival, with an argument for why
outcome-linked payment is expected to corrupt the underlying data.
\item A pre-registered evaluation protocol with named metrics, baselines and
falsification criteria, published in advance of data collection, together with
an explicit statement of the ecological, legal and equity limitations of the
approach.
\item A structured transferability analysis against the closest existing
deployment, Ericsson's Connected Mangroves in Malaysia, identifying which of its
findings generalise to the Kenyan setting and which of its architectural choices
are artefacts of a sponsorship and labour model that Kenya does not share.
\end{enumerate}

We deliberately report no accuracy figures, no pilot outcomes and no cost
savings. None have been measured. Section~\ref{sec:eval} specifies how they
should be measured and what result would falsify the central claim.

\section{Background and Related Work}

\subsection{Remote Sensing of Mangrove Extent and Condition}

Satellite monitoring of mangrove extent is mature. Sentinel-2 provides
multispectral imagery at 10~m ground sample distance for the visible and
near-infrared bands with a revisit interval of roughly five days, and its time
series have been used to map mangrove extent, track disturbance and characterise
recovery trajectories~\cite{s2monitoring,rsreview}. Multi-index approaches
combining vegetation, water and soil indices with machine learning classifiers
have been applied to mangrove dynamics assessment using Sentinel-2
imagery~\cite{multiindex}.

Two limitations bear directly on our design. First, resolution. At 10~m, a
single pixel covers 100~m$^2$. This is adequate for detecting stand-level
clearing and for tracking canopy closure across a restoration block, and wholly
inadequate for observing whether an individual propagule planted three months
ago is alive. Any claim that satellite data verifies planting survival at plot
level is unsupportable. Second, cloud. Optical revisit frequency is nominal, not
effective; the Kenyan coast has substantial cloud cover, and the usable optical
record is considerably sparser than the five-day figure implies. Synthetic
aperture radar, which penetrates cloud, is the standard mitigation and is
treated here as future work rather than as an implemented capability.

The consequence is a division of labour rather than a choice: satellite data
supplies hectare-scale context and loss alerting, and ground observation
supplies plot-scale truth. Neither substitutes for the other.

\subsection{MRV Cost as the Scaling Constraint}

The MRV literature supports treating verification cost as a first-order design
variable. Vanderheyden et al.~\cite{vanderheyden2026} decompose MRV cost into
monitoring approach, project design parameters such as farm count and monitoring
frequency, and protocol requirements, and find monitoring to be the dominant
cost component with field sampling the most labour-intensive element. Their
finding that digitalisation can halve data-entry time establishes that process
automation yields real savings, while their conclusion that smallholder
scalability likely requires policy support indicates that automation alone may
be insufficient. We take both findings seriously: BlueProof is positioned as a
cost-reduction instrument whose sufficiency is an empirical question, not a
settled one.

\subsection{Vision-Language Models for Field Species Identification}
\label{sec:vlmrelated}

This subsection is the reason the system is designed the way it is.

A natural design for participatory monitoring is to let a general-purpose
vision-language model identify what is in a submitted photograph. Recent
evidence argues strongly against relying on that capability without guardrails.
Zhou et al.~\cite{zhou2026} evaluate general-purpose VLMs in the 2--8B parameter
range against a specialised model on a 96-species identification task. They
report three findings that constrain any system of this kind. Every model
evaluated, general-purpose or specialist, degrades sharply when moving from
curated photographs to field imagery, with domain gaps of 9.6 to 26.6 percentage
points. The specialised model outperforms the general-purpose VLMs by 33.2 to
59.2 percentage points. And the general-purpose models fabricate species names
at rates of 5.9 to 9.6 percent.

Notably, the authors characterise the degradation as reflecting general image
legibility rather than a failure of fine-grained discrimination as such. That
distinction is actionable. It suggests that a VLM may be more reliable when
asked whether an image is legible and what condition is visible than when asked
to commit to a species label, and that legibility assessment is itself a useful
product of the model.

Three design consequences follow, and we adopt all three.

\begin{itemize}
\item \textbf{Do not make the VLM the taxonomic authority.} Species is declared
by the monitor, who is standing at the plot and typically planted it, and the
model corroborates or disputes that declaration. A dispute is a signal for
review, not an automatic overwrite.
\item \textbf{Make abstention a first-class output.} A fabrication rate near ten
percent is unacceptable in a record that a carbon buyer will rely on. The system
must be able to decline.
\item \textbf{Leave room for a specialist model.} Where taxonomic precision is
required, the architecture should admit a domain-trained classifier of the kind
that outperformed general VLMs in~\cite{zhou2026}, with the VLM retained for the
open-ended judgements it handles well.
\end{itemize}

\subsection{Participatory Monitoring and the Expertise Bottleneck}

Participatory and citizen-science monitoring has an established record in
ecology, and the appeal in this setting is obvious: the people best positioned
to observe a creek-side plot weekly are the people who live beside it and
planted it. The recurring difficulty is not willingness or diligence but
credibility. A carbon standard requires evidence that an auditor will accept,
and the conventional route to that acceptance is expert presence, which is the
precise cost that~\cite{vanderheyden2026} identifies as dominant and as
regressive against small projects.

This produces a bottleneck with a specific shape. Observation is cheap and
abundant; attestation is expensive and scarce. Systems that respond by sending
experts to observe do not scale. Systems that respond by simply trusting
self-report do not clear the credibility bar. The interesting design space lies
between: mechanisms that let abundant lay observation carry an attestation
burden it could not carry alone, by adding independent corroboration and by
making the incentive structure of the reporter inspectable. BlueProof is an
attempt to occupy that space, and the two corroborating signals it adds are
machine assessment of the submitted image and satellite observation of the
surrounding canopy.

\subsection{Gap}

Existing work establishes that satellite monitoring of mangroves is feasible at
canopy scale, that MRV cost is a dominant and regressive constraint, and that
general VLMs are unreliable taxonomic oracles on field imagery. What is missing
is a system design that accepts all three findings simultaneously: one that uses
satellite data only for what it can support, uses a VLM only for what the
evidence says it can do, and couples verification to a payment mechanism whose
incentive structure does not corrupt the data being verified. BlueProof is a
proposal for such a system.

\section{Precedent and Transferability: Connected Mangroves}
\label{sec:malaysia}

The proposition that networked technology can improve mangrove restoration
outcomes is not speculative. It has been demonstrated at field scale in
Malaysia. Because that deployment is the closest existing precedent to the
present work, and because it differs from the Kenyan setting in ways that are
instructive rather than incidental, we treat it as a comparative case rather
than a citation.

\subsection{The Malaysian Deployment}

Ericsson launched the Connected Mangroves project in Selangor, Malaysia, in
October 2015, working with a local technology partner, the Global Environment
Centre and the community of Kampung Dato Hormat~\cite{unfccc2015,ericsson2019}.
Sensor systems deployed among planted saplings measure soil moisture,
temperature and electrical conductivity together with ambient humidity and water
level, relay readings over commercial cellular networks, and present them
through a cloud dashboard accessible to community leaders, non-governmental
organisations, analysts and authorities, with automated alerts when thresholds
are breached~\cite{unfccc2015}. Each sensor installation covers roughly
2{,}500~m$^2$, on the order of 200 mangrove plants~\cite{unfccc2015}.

The reported results are meaningful. Against a regional baseline in which only
about 40 percent of planted seedlings reach maturity, the project targeted a 50
percent improvement in survival, beginning with 200 seedlings in a first phase
and scaling to 1{,}000~\cite{unfccc2015}. Subsequent reporting describes
survival of approximately 80 percent across roughly 4{,}000 saplings planted in
Malaysia, with the model replicated to a site in the
Philippines~\cite{ericsson2019}.

\subsection{What the Precedent Establishes, and What It Does Not}

Connected Mangroves establishes three things that we adopt without reservation.
Restoration survival is responsive to information: knowing that a plot is
drying or inundated, in time to act, materially changes outcomes. Community
custodianship is compatible with instrumentation. And a mangrove restoration
technology programme can be sustained for the better part of a decade, which is
longer than most pilots survive.

It does not, however, establish what BlueProof needs to establish, and the
difference is not one of degree. Connected Mangroves is an \emph{agronomic
telemetry} system. It measures the abiotic conditions in which seedlings grow,
so that people can intervene. It does not attest that a particular person
performed a particular act of stewardship at a particular place and time, which
is the artefact a carbon standard purchases and the artefact that community
projects cannot currently produce cheaply. A soil moisture reading tells you the
ground is dry. It does not tell you that Amina planted forty propagules on
Tuesday, nor does it pay her for having done so.

The two systems therefore solve adjacent but distinct problems: one optimises
growing conditions, the other manufactures verifiable evidence and couples it to
compensation.

\subsection{Points of Divergence}

Table~\ref{tab:malaysia} sets out the structural differences that bear on
transferability.

\begin{table*}[t]
\centering
\small
\caption{Structural comparison between the Malaysian Connected Mangroves
deployment and the constraints facing a Kenyan community mangrove project.}
\label{tab:malaysia}
\begin{tabular}{@{}p{3.1cm}p{6.2cm}p{6.2cm}@{}}
\toprule
& \textbf{Connected Mangroves, Malaysia} & \textbf{Kenyan coastal community project} \\
\midrule
Problem addressed & Seedling survival under variable conditions & Verifiable attestation for carbon finance \\
Primary sensing & Fixed abiotic sensors among saplings & Human observation of plot condition \\
Who bears the cost & Corporate sponsor supplies hardware and connectivity & Project must self-finance from verification value \\
Coverage density & One installation per {\raise.17ex\hbox{$\scriptstyle\sim$}}0.25~ha & Must cover tens of hectares per monitor \\
Community role & Volunteer planting and basic maintenance & Paid monitoring as a livelihood \\
Payment loop & None; information is the output & Conditional payment on verified submission \\
Governance context & State forestry with long silvicultural tradition & Gazetted reserve, co-managed, contested use \\
Infrastructure advantage & Cellular IoT backhaul & Near-universal mobile money \\
Dominant failure risk & Sensor maintenance and sponsor withdrawal & Verification credibility and model error \\
\bottomrule
\end{tabular}
\end{table*}

Four of these deserve elaboration.

\textbf{The sponsorship question is decisive.} Connected Mangroves is
underwritten by a telecommunications vendor supplying hardware, connectivity and
volunteer labour. That is a legitimate and generous arrangement, and it means
the project's unit economics are not the project's own. A Kenyan community
forest association has no standing equivalent. Any design that assumes donated
hardware inherits a dependency that, if withdrawn, ends the programme. BlueProof
therefore carries no per-site capital cost by construction, which is the
strongest single argument for community photography over instrumentation in
Table~\ref{tab:modalities}.

\textbf{Sensor density does not scale to a national estate.} At one installation
per 0.25~hectare, covering Kenya's 61{,}271 hectares~\cite{kairo2021} would
require on the order of $2.4 \times 10^{5}$ installations. Even the 42-hectare
site considered in this work would require roughly 170. The Malaysian density is
appropriate for an intensively studied demonstration plot and is not a national
monitoring architecture. Human observers scale differently: a monitor already
lives beside the creek, walks it anyway, and can cover hectares rather than
fractions of one.

\textbf{Kenya holds an infrastructure advantage the Malaysian project could not
use.} Mobile money is close to universal in Kenya, which makes it possible to
settle a payment to an individual monitor within seconds of a verification
passing. Connected Mangroves informs its stakeholders; it does not pay them per
verified act, and in a volunteer model there is no reason it should. The
conditional payment loop is not a refinement of the Malaysian design. It is a
mechanism Kenya can build that Malaysia's infrastructure and labour model did
not call for, and it converts conservation from an activity that costs a
community time into one that returns income.

\textbf{The Malaysian success metric is a trap if imported directly.} The
headline result is survival rising from roughly 40 percent to roughly 80
percent~\cite{unfccc2015,ericsson2019}. As a measure of ecological outcome that
is exactly right. As a basis for paying people it is hazardous, for the reason
set out in Section~\ref{sec:incentive}: once survival determines income, the
reporter acquires an interest in the number. Connected Mangroves is insulated
from this because its survival figures come from instruments rather than from
paid reporters. A Kenyan system that pays monitors and simultaneously grades
them on survival would not be insulated, and would corrupt the very statistic it
was built to produce.

\subsection{Implication for the Present Design}

The transferable lesson from Malaysia is the premise, not the apparatus.
Information improves restoration outcomes, and communities will sustain a
technology programme over years. The apparatus does not transfer, because the
Kenyan binding constraints are attestation cost, absence of a sponsor and the
need for the work to pay, none of which fixed environmental sensors address.

BlueProof can therefore be read as the Malaysian premise reconstructed under
Kenyan constraints: observation performed by people rather than instruments,
because people are already present and instruments must be bought and
maintained; verification performed by a model with a bounded and auditable role,
because expert attestation is the scarce resource; and compensation settled
immediately over mobile money, because that is the infrastructure Kenya actually
has and because a conservation system that pays the people doing the
conservation is more likely to persist than one that does not.

\section{System Design}

\subsection{Design Principles}

\textbf{P1: Separate the deterministic from the probabilistic.} Vegetation
index computation, cloud screening, confidence thresholds, payment arithmetic
and ledger construction are ordinary code with defined behaviour. Image
interpretation and natural language generation are probabilistic. The boundary
is explicit in the module structure so that each side can be audited on its own
terms. A recurring failure in applied machine learning for conservation is
delegating arithmetic or geometry to a language model; we avoid it by
construction.

\textbf{P2: Abstain rather than guess.} Given the fabrication rates reported
in~\cite{zhou2026}, the system treats low confidence as a routing decision, not
a discount factor. Submissions below a configured confidence threshold, or
whose assessed condition is \emph{unclear}, are escalated to human review and
trigger no payment.

\textbf{P3: Pay for verified reporting, not for outcomes.} This is discussed in
Section~\ref{sec:incentive}.

\textbf{P4: Degrade gracefully.} Field monitors work in locations with
intermittent connectivity and on inexpensive handsets. Every external dependency
in the reference implementation has a deterministic offline mode, and the
assistant falls back to rule-based responses if the model endpoint is
unavailable, so that a monitor in the field is never left without an answer.

\subsection{Architecture}

Figure~\ref{fig:pipeline} shows the chain of custody from field observation to
verified record.

\begin{figure}[t]
\centering
\small
\begin{tabular}{@{}c@{}}
\hline
\rule{0pt}{2.4ex}Monitor captures geotagged plot photograph\\
\hline
$\downarrow$\\
\hline
\rule{0pt}{2.4ex}VLM bounded assessment: legibility, condition,\\
species corroboration, cutting, pests, confidence\\
\hline
$\downarrow$\\
\hline
\rule{0pt}{2.4ex}Threshold gate: confidence $\geq \tau$ and condition known?\\
\emph{No} $\rightarrow$ human review queue, no payment\\
\hline
$\downarrow$ \emph{Yes}\\
\hline
\rule{0pt}{2.4ex}Conditional mobile-money payment to monitor\\
\hline
$\downarrow$\\
\hline
\rule{0pt}{2.4ex}Append-only stewardship ledger entry\\
\hline
$\downarrow$\\
\hline
\rule{0pt}{2.4ex}Satellite corroboration at site scale (NDVI, cloud-screened)\\
\hline
$\downarrow$\\
\hline
\rule{0pt}{2.4ex}Report generation: facts assembled in code,\\
narrative drafted by language model\\
\hline
\end{tabular}
\caption{The BlueProof verification and payment pipeline. Probabilistic
components are confined to the second and final stages; all gating arithmetic
is deterministic.}
\label{fig:pipeline}
\end{figure}

\subsection{Data Model and Chain of Custody}

The schema encodes provenance explicitly. A \emph{Site} carries a boundary
polygon, an area and a baseline vegetation index. A \emph{Plot} belongs to a
site and carries a field code, a declared species, a planting date, a count of
seedlings planted and a coordinate. A \emph{MonitoringEvent} belongs to a plot
and a monitor and carries both the monitor's declaration and the model's
assessment as separate fields, so that agreement and disagreement remain
recoverable after the fact. A \emph{SatelliteObservation} belongs to a site and
carries an index value, a change against baseline, a cloud fraction and an alert
flag. A \emph{Payment} belongs to a monitoring event. A \emph{LedgerEntry} is
written only after a payment settles and denormalises the verified facts into an
append-only record: site, plot, species, count verified, condition, the named
monitor, the amount paid and the coordinate.

Two properties of this schema matter for verification integrity. First, the
monitor's claim and the model's assessment are never merged into a single
field, so a later auditor can compute agreement rates and detect drift. Second,
the ledger records the amount paid alongside the ecological facts, which makes
the distribution of benefit inspectable rather than inferred. We return to why
that matters in Section~\ref{sec:equity}.

\subsection{Bounded Role of the Vision-Language Model}

The model is asked four questions about a submitted photograph, in this order of
confidence in its answers:

\begin{enumerate}[label=(\alph*)]
\item \textbf{Is this image legible enough to assess?} Following~\cite{zhou2026},
legibility is treated as the primary determinant of downstream reliability and
is therefore assessed explicitly rather than implied by a confidence score.
\item \textbf{What condition is visible?} Healthy, stressed, dead or unclear.
This is a coarse ordinal judgement over a visually salient property.
\item \textbf{Is there evidence of cutting or pest damage?} Stumps, cut poles and
axe marks; boring holes, frass and dieback. These are presence detections of
conspicuous features, and are treated as recall-oriented: a false positive costs
a human review, a false negative loses a deforestation event.
\item \textbf{Is the monitor's declared species consistent with the image?}
Corroboration, not determination. Disagreement routes to review.
\end{enumerate}

The model is given the nine mangrove species recorded in Kenya as a closed
label set, which reduces the space in which fabrication can occur, and any
returned label outside that set is coerced to \emph{unclear} rather than
accepted. The architecture reserves a slot for a domain-trained classifier to
replace the model on question~(d), consistent with the specialist advantage
reported in~\cite{zhou2026}.

\subsection{Satellite Corroboration}

For each site, a normalised difference vegetation index series is computed from
Sentinel-2 red and near-infrared bands over the site polygon and compared
against a stored baseline. Observations whose cloud fraction exceeds a
configured maximum are excluded from alerting rather than down-weighted, on the
principle that an untrustworthy observation should not contribute to a decision
that moves money or triggers an enforcement response. A sustained negative
change beyond a configured threshold raises a canopy-loss alert for the site.

The satellite layer is explicitly \emph{corroborative}. It cannot confirm plot
survival, and the system does not claim that it does. Its value is threefold: it
detects loss the field record might miss because nobody was sent to look, it
provides an independent signal against which systematic misreporting would show
up as divergence, and it supplies the landscape-scale context a report requires.

\subsection{Alternative Sensing Modalities Considered}

Community smartphone photography is not the only way to gather evidence from a
mangrove stand, and the choice deserves justification rather than assumption.
Table~\ref{tab:modalities} compares four candidate modalities against the
constraints that actually bind a small project in this setting.

\begin{table*}[t]
\centering
\small
\caption{Sensing modalities for community mangrove monitoring, assessed against
constraints binding on a small coastal project.}
\label{tab:modalities}
\begin{tabular}{@{}lllll@{}}
\toprule
& \textbf{Community photography} & \textbf{UAV survey} & \textbf{Passive acoustics} & \textbf{Fixed IoT sensors} \\
\midrule
Capital cost per site & None; existing handsets & High & Moderate per node & Moderate to high \\
Regulatory burden & None & Pilot licence, operator cert. & None & None \\
Realistic time to deploy & Days & Weeks to months & Weeks & Weeks \\
Field maintenance & None beyond the handset & Skilled operator & Retrieval, batteries & Power, theft, corrosion \\
Primary observable & Plot condition, disturbance & Canopy structure, extent & Faunal activity & Abiotic conditions \\
Plot-level survival & Yes & Partially & No & No \\
Community acceptance & High; participants are paid & Contested; surveillance framing & Neutral & Neutral \\
Generates livelihood & Yes & No & No & No \\
\bottomrule
\end{tabular}
\end{table*}

Two of these deserve comment. Uncrewed aerial survey produces the most
compelling imagery and is the weakest practical fit. In Kenya, drone operation
requires registration supported by police clearance, and operation beyond the
most restricted category requires a remote pilot licence together with
operational authorisation, with commercial operation requiring an operator
certificate; processing timelines are on the order of a
month~\cite{kcaa}. For a small project this converts a technical question into an
administrative one, and an autonomous patrol falls squarely in the categories
requiring prior authorisation. There is also a social cost that is easy to
overlook: aerial cameras operating over a settlement read as surveillance
regardless of intent, which is corrosive in a context where the system depends on
voluntary participation.

Passive acoustic monitoring is scientifically attractive, particularly for
faunal recovery, which is a genuine gap in vegetation-only assessment. It does
not observe seedling survival, requires hardware deployment and retrieval, and
depends on labelled regional audio that does not presently exist for this
coastline. We regard it as a valuable complement at a later stage rather than a
foundation.

The decisive argument for community photography is not sensing performance. It
is that it is the only modality in Table~\ref{tab:modalities} that transfers
income to the community performing the conservation work. The observation and
the livelihood are the same act. Every other modality treats the community as a
subject of measurement rather than a participant in it.

\subsection{Conditional Payment}
\label{sec:incentive}

Payment is released only after a monitoring event passes the verification gate,
and is delivered over mobile money, which in the Kenyan context means a
near-instant transfer to a handset the monitor already owns.

The design decision that requires defending is what the payment is \emph{for}.
BlueProof pays a flat amount per verified submission. It does not pay per
surviving seedling.

The argument is one of incentive compatibility. If payment is proportional to
survival, a monitor who discovers that half a plot has died faces a direct
financial penalty for reporting it accurately. The rational response is to
photograph the surviving half, to report optimistically, or to replant before
the assessment. Every one of those responses degrades exactly the data that the
carbon buyer is purchasing, and the degradation is invisible precisely because
the monitoring system has been trained on the resulting data. Paying for honest,
verifiable reporting inverts this: mortality reported is still paid work, so the
monitor's interest and the buyer's interest coincide. The verified record
becomes worth something because it is not distorted by the mechanism that
produced it.

This decision has a cost. It decouples payment from ecological outcome, which
some funders will dislike, and it means that a monitor can in principle be paid
for reporting a failed restoration. We consider that the correct trade: the
alternative purchases a number that looks better and means less. We propose a
field experiment to test the claim empirically in Section~\ref{sec:eval}.

\subsection{Report Generation}

Monitoring standards require a structured narrative. In BlueProof the
quantitative content of that narrative is assembled deterministically from the
ledger and the satellite series, and the language model receives those figures
as structured input with instructions to use only the supplied numbers, to state
plainly where evidence is weak or contradictory, and to flag divergence between
the field record and the satellite signal. The model is never asked to retrieve
or estimate a quantity. If the model endpoint is unavailable, a deterministic
template produces the same figures without narrative prose, so the report
pipeline has no hard dependency on model availability.

\section{Reference Implementation}

A reference implementation exists and exercises the full pipeline end to end.
The backend is a Python service built on FastAPI with SQLModel over SQLite,
chosen so that the system runs with no infrastructure provisioning during
evaluation; the storage layer is portable to PostgreSQL, and the geometry
representation to PostGIS, when spatial querying becomes necessary. The service
exposes endpoints for site and plot registration, photograph assessment,
monitoring submission with integrated verification and payment, satellite series
refresh, report generation and an impact rollup, together with a monitor-facing
assistant endpoint intended for a messaging channel.

Table~\ref{tab:api} summarises the service surface and, for each operation,
which side of the deterministic--probabilistic boundary it falls on.

\begin{table}[t]
\centering
\small
\caption{Service operations and their computational character. Probabilistic
components are confined to three operations, all of which have deterministic
fallbacks.}
\label{tab:api}
\begin{tabular}{@{}p{3.35cm}p{4.0cm}@{}}
\toprule
\textbf{Operation} & \textbf{Character} \\
\midrule
Register site, plot & Deterministic \\
Assess photograph & Probabilistic; closed label set \\
Submit monitoring event & Deterministic gate over a probabilistic assessment \\
Release payment & Deterministic; gated \\
Write ledger entry & Deterministic; append-only \\
Refresh satellite series & Deterministic; cloud-screened \\
Generate report & Facts deterministic; narrative probabilistic \\
Field assistant & Probabilistic; rule-based fallback \\
Impact rollup & Deterministic \\
\bottomrule
\end{tabular}
\end{table}

Every external dependency has a deterministic offline mode: the vision model,
the satellite source and the payment rail can each be independently switched to
a live provider through configuration without code change. This is a testing
convenience, but it is also a field requirement, since demonstrations and
training sessions frequently happen without reliable connectivity.

The monitor-facing assistant answers practical questions in Swahili with some
English, covering species zonation, spacing, how to record a plot and when
payment is released. Zonation guidance follows the standard structure of Kenyan
mangrove stands, with \emph{Rhizophora mucronata} in regularly inundated lower
zones, \emph{Avicennia marina} in higher and more saline areas, \emph{Ceriops
tagal} intermediate, and \emph{Sonneratia alba} seaward.

We report no performance measurements for this implementation. It has been
exercised only against synthetic inputs.

\section{Proposed Evaluation}
\label{sec:eval}

Because no field data has been collected, we specify the evaluation in advance,
including the results that would falsify the central claim.

\subsection{Research Questions}

\textbf{RQ1.} Can a general-purpose VLM, restricted to a closed set of nine
Kenyan mangrove species, corroborate a monitor-declared species on field
photographs at a rate useful for triage, and at what abstention rate?

\textbf{RQ2.} What is the model's recall for conspicuous disturbance, namely
cutting and pest damage, relative to an expert field assessment?

\textbf{RQ3.} Does the verification pipeline reduce MRV cost per hectare per
monitoring cycle relative to current practice in a comparable community project?

\textbf{RQ4.} Does paying for verified reporting, rather than for survival,
reduce misreporting of seedling mortality?

\subsection{Dataset}

We propose a corpus of approximately 1{,}000 geotagged plot photographs
collected at Tudor Creek in Mombasa County and at Gazi Bay in Kwale County, the
latter chosen because an established research baseline exists there. Images
should be captured by community monitors on their own handsets, under the
lighting and weather conditions that will obtain in deployment, rather than by
researchers under favourable conditions; the domain gap reported
in~\cite{zhou2026} is precisely the effect that a curated corpus would hide.
Each image requires an independent expert label for species, condition and
presence of disturbance, established by a trained botanist in the field at time
of capture. Inter-rater agreement should be reported on a duplicated subset.

\subsection{Metrics}

\begin{itemize}
\item \textbf{Corroboration agreement}: rate at which the model's species
assessment agrees with the expert label, conditioned on the monitor's
declaration being correct and separately on it being incorrect. The second case
is the operationally important one, since it measures whether the system catches
a mistaken declaration.
\item \textbf{Disturbance recall and precision}, with recall privileged. We
propose reporting the precision obtained at a recall of 0.90 for cutting
detection, since a missed clearing event is a materially worse error than an
unnecessary human review.
\item \textbf{Abstention rate and abstention quality}: the fraction of
submissions routed to human review, and the error rate within the accepted set
compared with the error rate within the abstained set. An abstention mechanism
that does not concentrate errors in the abstained set is not working.
\item \textbf{Fabrication rate}: frequency of species labels outside the closed
set, or of asserted features absent from the image, measured against expert
labels. This is a direct replication target for the 5.9 to 9.6 percent range
reported in~\cite{zhou2026} under a constrained label set.
\item \textbf{Calibration}: expected calibration error of the reported
confidence against observed correctness, since the confidence threshold is the
mechanism that gates payment and must be meaningful.
\end{itemize}

\subsection{Baselines}

Four comparators are required for the result to be interpretable: the
monitor's unaided declaration; a domain-trained specialist classifier of the
kind shown to outperform general VLMs in~\cite{zhou2026}; an expert field
officer, representing current practice and the cost baseline; and a trivial
majority-class predictor, to establish the floor.

\subsection{Cost Study}

RQ3 requires an economic measurement, not a software benchmark. We propose
recording, for one monitoring cycle at a comparable community project, the
labour hours, travel, expert time and administrative overhead currently consumed
per hectare, and comparing against the same quantities under the proposed
workflow, inclusive of human review time generated by abstention and of model
inference cost. Abstention is a cost transfer, not a cost elimination, and an
evaluation that omits review labour would overstate the saving.

\subsection{Incentive Experiment}

RQ4 requires a field experiment with human participants and therefore ethical
review, informed consent in the participants' language, and the agreement of the
relevant community forest association. We propose assigning monitoring groups to
payment-for-reporting or payment-for-survival conditions and comparing reported
mortality against independently assessed mortality on a sample of plots. The
prediction is that reported mortality under payment-for-survival is
significantly lower than independently assessed mortality, and that the gap is
smaller under payment-for-reporting.

\subsection{Threats to Validity}

\textbf{Construct validity.} Expert labels are treated as ground truth, but
species determination from a photograph is genuinely difficult for young
individuals of morphologically similar taxa, and an expert working from the same
image rather than from the plant in hand inherits the same limitation. We
require labels established in the field at time of capture, and inter-rater
agreement reported, so that the ceiling imposed by label noise is visible rather
than assumed away.

\textbf{Internal validity.} Monitors who know their submissions are being
evaluated may photograph more carefully than they would in routine operation,
inflating measured legibility. A period of routine operation before the
evaluation window, and sampling of submissions without prior notice, mitigate
this only partially.

\textbf{External validity.} Two sites on one coastline, with one dominant set of
species and one set of lighting and tidal conditions, will not characterise
performance elsewhere. Results should be read as specific to Kenyan coastal
mangrove and not generalised to other mangrove regions, still less to other
ecosystems.

\textbf{Temporal validity.} Model behaviour changes as providers update
endpoints. Any reported figure is a measurement of one model version at one
point in time, and the evaluation should record the exact model identifier and
be treated as requiring periodic repetition rather than as a durable property of
the system.

\textbf{Conflict of interest.} The evaluation of a cost-reduction claim is being
proposed by the party who benefits from that claim being true. Cost measurement
in particular should be performed or audited by the comparison project rather
than by the system's author.

\subsection{Falsification Criteria}

The central claim of this paper fails if any of the following hold. If
corroboration agreement on genuinely mistaken declarations is no better than
chance, the verification layer adds no value over self-report. If the abstention
rate required to reach an acceptable error rate exceeds the fraction of
submissions a human reviewer could process at the target scale, the system does
not reduce expert dependence and the cost argument collapses. If the measured
MRV cost per hectare is not materially below current practice once review labour
is included, the premise of the work is wrong. We commit to reporting these
outcomes if they obtain.

\section{Discussion}

\subsection{What This System Does Not Do}

BlueProof is not a carbon registry and does not issue credits. It produces
evidence intended for use within an existing accredited methodology. Any
sequestration figure derived from its outputs depends on a per-hectare rate that
must come from an approved methodology or from site-specific measurement of the
kind reported in~\cite{kairo2021}, not from the system itself. The reference
implementation carries such a rate as an explicitly labelled configuration
assumption for demonstration purposes, and we would regard quoting a figure
derived from it as a misrepresentation.

\subsection{Equity and Benefit Sharing}
\label{sec:equity}

There is a legitimate criticism of blue carbon finance that a limited share of
revenue reaches the people performing the physical work. A system that reduces
verification cost does not address this, and could plausibly worsen it: cheaper
verification expands the supply of credits without altering how the proceeds are
divided, and the efficiency gain may simply accrue to intermediaries.

We regard this as a design obligation rather than an externality. Two mitigations
are built into the proposal. Recording the amount paid to a named monitor in the
same ledger entry as the ecological facts makes the labour share of each verified
record computable rather than obscured, and publishing the payment per verified
event as a project parameter makes it contestable. Neither guarantees equitable
distribution. Both make inequitable distribution visible, which is a
precondition for contesting it. We would encourage any deployment to publish the
ratio of monitor payments to total project revenue as a standing metric.

\subsection{Where the Cost Actually Moves}

It is worth being precise about the mechanism by which this design could reduce
cost, because a vague claim of automation invites the reasonable objection
that machine assessment simply relocates work.

Current practice concentrates cost in expert field presence: travel to a site,
time on the ground, and the administrative overhead of turning field notes into
an auditable narrative. The proposed workflow leaves observation with people who
are already present, which removes travel almost entirely for routine cycles,
and replaces a portion of expert attestation with machine corroboration plus
human review of the residual. Report preparation, which~\cite{vanderheyden2026}
associates with substantial administrative burden and where digitalisation is
shown to yield measurable savings, is largely automated because the quantitative
content is already structured.

Three components of cost do not reduce and must appear in any honest accounting.
Human review of abstained submissions is real labour, and its volume is set by
the abstention rate, which is in turn set by how conservative the confidence
threshold must be to hold the error rate acceptable. Model inference has a
per-submission cost. And initial training of monitors, together with periodic
re-training, is unavoidable and is not a one-off. The cost study specified in
Section~\ref{sec:eval} is constructed to capture all three, because an
evaluation that counted only the removed travel would produce a flattering and
false number.

\subsection{Governance and Tenure}

Mangroves in Kenya are gazetted government forest reserve. Any field deployment
requires the participation of the Kenya Forest Service and a community forest
association, and a system that generates a record of who did what where in a
contested resource carries governance implications that software design cannot
resolve. Photographs taken in and around a settlement are personal data in
practice regardless of their classification in law, and retention, access and
consent should be settled before collection begins, not after.

\section{Limitations}

We summarise the limitations stated throughout. The verification premise is
untested on real mangrove photography, and the published evidence gives
substantial reason for caution~\cite{zhou2026}. Satellite corroboration is
implemented against a synthetic index series; live ingestion is specified but
not built, and radar fallback for cloud is future work. The satellite layer
cannot observe plot-level survival at 10~m resolution, and 70.3 percent of the
carbon at stake in this ecosystem is in soil that no optical sensor
observes~\cite{kairo2021}. The cost thesis is argued from the general MRV
literature~\cite{vanderheyden2026} rather than measured in this setting. The
incentive argument is theoretical pending the experiment proposed above. And the
system has been exercised only against synthetic inputs, so no claim about its
behaviour under real load or real imagery is made here.

\section{Conclusion}

Kenya's mangroves store carbon at exceptional density, are being lost at a
measurable rate, and are protected in places by community institutions that
work but do not scale. We have argued that verification cost is the plausible
binding constraint, and presented a system design that attacks it while
respecting two constraints the evidence imposes: that satellite data cannot see
what people claim it sees, and that general-purpose vision-language models are
unreliable taxonomic authorities on field imagery. The resulting design confines
the model to bounded judgements, requires it to abstain, keeps monitor claims
and model assessments separately recorded, and pays for honest reporting rather
than for favourable outcomes.

Whether this reduces MRV cost enough to matter is an empirical question that
this paper does not answer. We have specified how to answer it, and what answer
would tell us the idea is wrong.

\section*{Acknowledgements}

This work is at design stage and has no institutional sponsor. No community
forest association, research institute or government agency has endorsed it, and
none should be understood to have done so. The author thanks the reviewers of
earlier drafts.

\begin{thebibliography}{9}
\small

\bibitem{kairo2021}
J. Kairo, A. Mbatha, M. M. Murithi, and F. Mungai,
``Total ecosystem carbon stocks of mangroves in Lamu, Kenya, and their potential
contributions to the climate change agenda in the country,''
\emph{Frontiers in Forests and Global Change}, vol. 4, 2021.
doi:10.3389/ffgc.2021.709227

\bibitem{zhou2026}
W. Zhou, M. Siripuram, X. Yan, Z. Liu, and Y. Ding,
``Can edge-deployable vision-language models identify species?''
\emph{arXiv preprint} arXiv:2609.11916, 2026.

\bibitem{vanderheyden2026}
L. Vanderheyden, N. Gayot, and G. Van den Broeck,
``What monitoring, reporting and verification (MRV) systems can reduce costs and
enhance scalability of carbon farming?''
\emph{Journal of Environmental Management}, vol. 401, 2026.

\bibitem{planvivo}
Plan Vivo Foundation,
``Mikoko Pamoja, Kenya,'' project documentation.
\url{https://www.planvivo.org/projects/mikoko-pamoja-kenya}

\bibitem{s2monitoring}
``Mangrove monitoring maps: tracking status, disturbance, and recovery
trajectories from Sentinel-2 time series,''
\emph{Science of Remote Sensing}, 2026.

\bibitem{rsreview}
``Mangrove mapping and monitoring using remote sensing techniques towards
climate change resilience,''
\emph{Scientific Reports}, vol. 14, 2024.
doi:10.1038/s41598-024-57563-4

\bibitem{unfccc2015}
United Nations Framework Convention on Climate Change,
``Connected Mangroves, Malaysia,'' Momentum for Change, ICT Solutions.
\url{https://unfccc.int/climate-action/momentum-for-change/ict-solutions/connected-mangroves}

\bibitem{ericsson2019}
Ericsson,
``IoT environmental monitoring for vital predictions: revitalizing mangrove
forests,'' case study, 2019.
\url{https://www.ericsson.com/en/cases/2019/revitalizing-mangrove-forests}

\bibitem{matang}
``Is Matang mangrove forest in Malaysia sustainably rejuvenating after more than
a century of conservation and harvesting management?''
\emph{PLOS ONE}, 2014. doi:10.1371/journal.pone.0105069

\bibitem{leemalaysia}
S. L. Lee, S. Y. Chee, M. Huxham, M. Jamilah, J. Choo, C. R. Kaur,
A. A. Amir, J. L. S. Ooi, M. Rozaimi, H. Omar, S. Sharma, M. Moritz,
and A. Y.-H. Then,
``Blue carbon ecosystems in Malaysia: status, threats, and the way forward for
research and policy,''
\emph{Journal of Environment and Development}, 2025.
doi:10.1177/10704965241284366

\bibitem{kcaa}
Kenya Civil Aviation Authority,
``Unmanned aircraft systems: remote pilot licensing and operator certification,''
regulatory guidance.
\url{https://www.kcaa.or.ke/safety-security-oversight/unmanned-aircraft-systems}

\bibitem{multiindex}
``Integrating multi-index remote sensing and machine learning for mangrove
dynamics assessment using Sentinel-2 imagery,''
\emph{iScience}, 2026.

\end{thebibliography}

\end{document}
